\documentclass[11pt,twocolumn]{article}
\usepackage[utf8]{inputenc}
\usepackage[T1]{fontenc}
\usepackage{times}
\usepackage{geometry}
\geometry{a4paper, margin=2.5cm}
\usepackage{amsmath,amssymb}
\usepackage{booktabs}
\usepackage{multirow}
\usepackage{graphicx}
\usepackage{xcolor}
\usepackage{hyperref}
\usepackage{natbib}
\usepackage{tabularx}
\usepackage{array}
\usepackage{enumitem}

\setlength{\columnsep}{0.6cm}

\title{Differential Depth Response in Iterative LLM Generation:\\
A Distributional Fingerprint for Provenance Characterization}

\author{Anonymous}
\date{}

\begin{document}
\maketitle

%=============================================================================
% ABSTRACT
%=============================================================================
\begin{abstract}
Web-scraped training corpora increasingly contain LLM-generated text of unknown provenance depth.
Current detectors offer only binary real-vs-synthetic classification, conflating qualitatively different material:
first-generation text retains much of the source lexical distribution, while deeper recursive outputs exhibit measurably reduced vocabulary diversity and compressed perplexity.
We formulate \textbf{generation depth classification} (gen-0/gen-1/gen-2+) and characterize \textbf{differential depth response} (DDR) across 8 base-model generators from 5 families (0.6B--8B):
central-tendency features (entropy mean, TTR, hapax ratio) drift monotonically with depth ($|\rho| \geq 0.55$), while shape features (tail decay, entropy max) remain near-stationary ($|\rho| < 0.10$).
A lightweight 7-feature ensemble achieves leak-free binary F1 = 0.694 under 8-generator leave-one-generator-out (LOGO) evaluation (95\% CI [0.688, 0.700]), and 0.830 on cross-model chains.
Cross-family generalization is confirmed via leave-one-family-out evaluation (binary F1 = 0.686).
DDR is robust across 3 reference models from different families; shape-feature stationarity holds with $|\rho| < 0.06$.
Depth-aware data curation is statistically equivalent to random subsampling for downstream causal language modeling quality (TOST $p = 1.1 \times 10^{-4}$ at $\pm 2\%$ margin), establishing depth fingerprinting as a \textbf{diagnostic} tool for provenance characterization rather than a curative tool for data selection.
Current scope is restricted to base models at moderate temperatures ($t \leq 0.9$).
\end{abstract}

%=============================================================================
% 1. INTRODUCTION
%=============================================================================
\section{Introduction}

Large language models (LLMs) are generating an increasing fraction of web content.
Recent estimates suggest 7--17\% of scientific papers and significant portions of web-scraped corpora contain LLM-generated text \citep{liang2024monitoring}.
This poses challenges for training data curation: recursive consumption of synthetic data can degrade model quality through distributional collapse \citep{shumailov2024model}.

Current AI text detection methods operate as binary classifiers --- real vs.\ synthetic --- and cannot distinguish \emph{how many generations removed} a text is from its human-authored seed.
This distinction matters: first-generation (gen-1) text retains much of the source distribution's lexical properties, while deeper recursive outputs exhibit measurably reduced vocabulary diversity and compressed perplexity distributions.
Conflating these qualitatively different materials limits our ability to understand and mitigate model collapse risks.

We address this gap by formulating \textbf{generation depth classification}: given a text, predict whether it is human-authored (gen-0), first-generation LLM output (gen-1), or deeper recursive output (gen-2+).
Our investigation reveals \textbf{differential depth response} (DDR): features derived from the next-token log-likelihood (NLL) distribution exhibit a systematic split.
\emph{Central-tendency features} (entropy mean, type-token ratio, hapax ratio) drift monotonically with generation depth, consistent with recursive compression of the token distribution.
\emph{Shape features} (tail decay rate, entropy max) remain near-stationary, indicating that recursive generation induces a \emph{location shift} in the NLL distribution while approximately preserving distributional shape.

We validate DDR across 8 base-model generators from 5 architectural families (Pythia, OLMo, Qwen, Llama, Mistral; 0.6B--8B parameters) and 3 reference models.
A 7-feature ensemble achieves leak-free binary F1 = 0.694 under LOGO evaluation, outperforming repurposed binary detectors (avg-logprob 0.765, FastDetectGPT 0.627, Binoculars 0.579 on macro F1).
Cross-model chains --- where each generation step uses a different generator --- yield binary F1 = 0.830, suggesting generator heterogeneity \emph{amplifies} rather than obscures depth signal.

Our third contribution is a calibrated negative finding: depth-aware data curation does not improve downstream language modeling quality compared to random subsampling (TOST equivalence confirmed at $\pm 2\%$ margin).
This positions depth fingerprinting as a \textbf{diagnostic tool} for data provenance characterization, not a curative tool for data selection.

We report scope limitations transparently: the signal degrades substantially under instruction tuning (LOGO F1 = 0.338 vs.\ 0.763 for base models), restricting current applicability to base models at moderate sampling temperatures.

%=============================================================================
% 2. RELATED WORK
%=============================================================================
\section{Related Work}

\paragraph{AI text detection.}
Perplexity-based methods leverage distributional differences between human and machine text.
GLTR \citep{gehrmann2019gltr} visualizes token-level log probabilities;
DetectGPT \citep{mitchell2023detectgpt} and its faster variant Fast-DetectGPT \citep{bao2024fast} use perturbation-based log probability curvature;
Binoculars \citep{hans2024spotting} computes cross-entropy ratios between two models;
DNA-GPT \citep{yang2024dna} uses divergent n-gram analysis.
All operate as binary classifiers and cannot estimate generation depth.

\paragraph{Model collapse \& iterative generation.}
\citet{shumailov2024model} formalize model collapse under iterative training on synthetic data, showing distributional tails vanish over generations.
\citet{dohmatob2024tale} provide theoretical bounds on variance collapse.
LLM as Broken Telephone \citep{briakou2025llm} studies paraphrastic attractor cycles, showing convergence within $\sim$3 iterations.
Our work addresses the complementary non-paraphrastic setting where unconstrained generation permits unbounded distributional drift.

\paragraph{Data quality \& provenance.}
Vocabulary diversity metrics (TTR, hapax ratio) are established in corpus linguistics \citep{mccarthy2010mtld}.
LD-Score \citep{tulchinskii2024intrinsic} uses character-level distribution divergence as a detection signal.
Perplexity-based data filtering \citep{wenzek2020ccnet} is standard in corpus construction.
Our work extends these to ordinal depth estimation.

\paragraph{Stylometry \& authorship.}
Authorship attribution methods \citep{neal2017surveying} use distributional features for author identification.
Our feature set draws on this tradition but targets generation depth rather than individual authorship.

%=============================================================================
% 3. METHOD
%=============================================================================
\section{Method}

\subsection{Task Formulation}

Given a text passage $x$, we predict its \emph{generation depth} $d \in \{0, 1, 2+\}$:
\begin{itemize}[nosep,leftmargin=*]
    \item \textbf{gen-0}: Human-authored seed text.
    \item \textbf{gen-1}: First-generation LLM output from a human seed.
    \item \textbf{gen-2+}: Output of $\geq$2 recursive generation steps.
\end{itemize}

The \textbf{primary metric} is leak-free binary F1 (gen-1 vs.\ gen-2+), excluding gen-0 to avoid seed-sharing leakage.
Three-class macro F1 is reported but noted as inflated by gen-0's strong discriminability.

\subsection{Generation Chain Construction}

For each human-authored seed $s$, we construct a chain $\{x_0 = s, x_1, x_2, x_3, x_4\}$ where $x_{d+1} = \text{LLM}(x_d)$ using continuation prompting at temperature $t = 0.7$.
Chains are constructed independently per generator.
Cross-model chains use a different generator at each depth step.

\subsection{Feature Extraction}
\label{sec:features}

We extract 7 features per passage using a reference model $M_{\text{ref}}$ (Pythia-410M by default):

\paragraph{Surface features} (computed from raw text):
\begin{itemize}[nosep,leftmargin=*]
    \item \textbf{TTR}: Type-token ratio (unique tokens / total tokens).
    \item \textbf{Hapax ratio}: Fraction of tokens appearing exactly once.
\end{itemize}

\paragraph{Reference features} (computed from $M_{\text{ref}}$'s NLL distribution over $x$):
\begin{itemize}[nosep,leftmargin=*]
    \item \textbf{Entropy mean}: Mean per-token NLL.
    \item \textbf{Entropy std}: Standard deviation of per-token NLL.
    \item \textbf{Entropy min/max}: Extremal NLL values.
    \item \textbf{Tail decay}: Decay rate of the upper tail of the NLL distribution, estimated via exponential fit.
\end{itemize}

\subsection{Differential Depth Response (DDR)}

DDR characterizes feature behavior as a function of generation depth via Spearman rank correlation $\rho$ between feature value and depth $d$:
\begin{itemize}[nosep,leftmargin=*]
    \item \textbf{Drifting features} ($|\rho| \geq 0.55$): hapax ratio, TTR, entropy mean. These exhibit monotonic trends with depth, consistent with recursive compression.
    \item \textbf{Stationary features} ($|\rho| < 0.10$): tail decay, entropy max. These remain near-constant, indicating the NLL distribution's shape is preserved under location shift.
    \item \textbf{Weakly drifting}: entropy std ($|\rho| \approx 0.20$), intermediate between the two groups.
\end{itemize}

\subsection{Evaluation Protocol}

We use two leave-out evaluation protocols to prevent information leakage:

\paragraph{LOGO} (Leave-One-Generator-Out): Train on $K{-}1$ generators, test on the held-out generator. Measures cross-generator generalization.

\paragraph{LOFO} (Leave-One-Family-Out): Train on all generators except those from one architectural family (e.g., hold out all Pythia models). Measures cross-family transfer.

The classifier is XGBoost with default hyperparameters, evaluated via 5-fold cross-validation within each training split.

%=============================================================================
% 4. EXPERIMENTS
%=============================================================================
\section{Experiments}

\subsection{Setup}

\paragraph{Generators.} We evaluate 8 base-model generators spanning 5 families: Pythia-1B, Pythia-2.8B (EleutherAI); OLMo-2-1B (AI2); Qwen3-0.6B, Qwen3-1.7B, Qwen3-8B (Alibaba); Llama-3.2-3B (Meta); Mistral-7B (Mistral AI).
All generation uses continuation prompting at $t = 0.7$.

\paragraph{Seeds.} Wikipedia passages serve as the primary seed domain ($\sim$1000 seeds per generator), with additional evaluation on news, academic, and code domains (\S\ref{sec:multidomain}).

\paragraph{Reference model.} Pythia-410M (default). Cross-reference robustness evaluated with OLMo-2-1B and Qwen3-0.6B (\S\ref{sec:crossref}).

\subsection{Feature Ablation \& Trend Analysis}
\label{sec:ablation}

Table~\ref{tab:ablation} reports 5-fold CV results on the full dataset, ablating feature groups.

\begin{table}[t]
\centering
\small
\begin{tabular}{lc}
\toprule
\textbf{Feature Set} & \textbf{Macro F1} \\
\midrule
Full (7 features) & 0.820 \\
Surface only (TTR, hapax) & 0.763 \\
Reference only (5 features) & 0.719 \\
\bottomrule
\end{tabular}
\caption{Feature ablation (5-fold CV). Full set significantly outperforms both subsets (McNemar $p < 0.001$), confirming surface--reference complementarity.}
\label{tab:ablation}
\end{table}

Surface and reference features are complementary: the full set outperforms surface-only by 5.7 pp and reference-only by 10.1 pp (McNemar $p < 0.001$ for both).
This establishes that depth signal arises from two distinct sources: vocabulary diversity reduction (surface) and NLL distribution shift (reference).

\paragraph{DDR trend analysis.}
Linear regression of features against depth confirms the drifting/stationary split (Table~\ref{tab:trends}).
Drifting features show high $R^2$: entropy mean (0.93), hapax ratio (0.92), TTR (0.89).
Stationary features show near-zero $R^2$: tail decay (0.03).
Entropy std occupies an intermediate position ($R^2 = 0.41$, slope $= -0.12$), with drift magnitude $\sim$1/3 of the drifting group.

\begin{table}[t]
\centering
\small
\begin{tabular}{lccc}
\toprule
\textbf{Feature} & \textbf{Slope} & $R^2$ & \textbf{Category} \\
\midrule
Entropy mean & $-0.29$ & 0.93 & Drifting \\
Hapax ratio & $-0.043$ & 0.92 & Drifting \\
TTR & $-0.034$ & 0.89 & Drifting \\
Entropy std & $-0.12$ & 0.41 & Weak \\
Tail decay & $+0.002$ & 0.03 & Stationary \\
\bottomrule
\end{tabular}
\caption{Per-feature linear trend across depths 0--4. Drifting features exhibit strong monotonic decrease ($R^2 > 0.89$); tail decay is effectively flat.}
\label{tab:trends}
\end{table}

Perplexity under the reference model decreases monotonically with depth (median: 13.5, 4.96, 3.00, 2.07, 1.72 for depths 0--4), confirming that recursive generation compresses the NLL distribution leftward.

\subsection{Cross-Generator Generalization}
\label{sec:crossgen}

Table~\ref{tab:crossgen} reports per-generator LOGO results on the 5-generator subset (Pythia-2.8B, Pythia-1B, OLMo-2-1B, Qwen3-1.7B, Qwen3-0.6B; $N = 23{,}032$).

\begin{table}[t]
\centering
\small
\begin{tabular}{lcc}
\toprule
\textbf{Generator} & \textbf{Macro F1} & \textbf{Binary F1} \\
\midrule
Pythia-1B & 0.851 & 0.724 \\
Pythia-2.8B & 0.804 & 0.720 \\
Qwen3-0.6B & 0.818 & 0.715 \\
Qwen3-1.7B & 0.805 & 0.710 \\
OLMo-2-1B & 0.752 & 0.670 \\
\midrule
\textbf{Average} & \textbf{0.804} & \textbf{0.720} \\
\bottomrule
\end{tabular}
\caption{Per-generator LOGO evaluation (5-gen subset). Binary F1 variance across generators is 5.4 pp, confirming cross-generator robustness.}
\label{tab:crossgen}
\end{table}

The LOGO binary F1 of 0.720 (95\% CI [0.708, 0.732]) demonstrates cross-generator generalization.
Per-generator variance is 5.4 pp, with OLMo-2-1B lowest (0.670) due to its non-GPT-NeoX architecture producing more divergent distributional patterns.
Hapax ratio is the most important feature (importance = 0.289), followed by TTR (0.168) and entropy mean (0.160).

\paragraph{Cross-family transfer.}
LOFO evaluation confirms generalization beyond same-family models: Pythia held-out F1 = 0.690, OLMo held-out = 0.720, Qwen held-out = 0.706 (average LOFO binary F1 = 0.705).
The 1.5 pp drop from LOGO to LOFO indicates minimal family-specific inflation.

\subsection{Cross-Reference-Model Robustness}
\label{sec:crossref}

To verify that DDR is not an artifact of the Pythia-410M reference model, we extract features using two additional reference models: OLMo-2-1B and Qwen3-0.6B (Table~\ref{tab:crossref}).

\begin{table}[t]
\centering
\small
\begin{tabular}{lccc}
\toprule
\textbf{Reference} & \textbf{Macro F1} & \multicolumn{2}{c}{\textbf{Spearman $\rho$}} \\
\cmidrule(lr){3-4}
 & & hapax & tail decay \\
\midrule
Pythia-410M & 0.803 & $-0.647$ & $-0.040$ \\
OLMo-2-1B & 0.807 & $-0.647$ & $-0.038$ \\
Qwen3-0.6B & 0.816 & $-0.647$ & $-0.042$ \\
\bottomrule
\end{tabular}
\caption{LOGO macro F1 and representative Spearman $\rho$ by reference model. DDR pattern (drifting vs.\ stationary split) is consistent across all three reference families. Surface features are identical by construction.}
\label{tab:crossref}
\end{table}

Classification F1 varies by only 1.3 pp across reference models.
Entropy mean drift direction is consistently negative across all three reference families (Pythia: $[-1.19, -1.40]$; OLMo: $[-0.61, -0.91]$; Qwen: $[-0.96, -1.35]$).
Shape-feature stationarity holds robustly: maximum $|\rho|$ across all reference models is 0.055.
DDR is a property of the generation process, not an artifact of reference model choice.

\subsection{Baseline Comparison}
\label{sec:baseline}

We compare our 7-feature ensemble against four existing detectors repurposed as ordinal depth predictors (Table~\ref{tab:baseline}).

\begin{table}[t]
\centering
\small
\begin{tabular}{lcc}
\toprule
\textbf{Method} & \textbf{LOGO F1} & \textbf{Spearman $\rho$} \\
\midrule
7-feature ensemble & \textbf{0.798} & --- \\
Avg log-prob & 0.765 & 0.788 \\
LD-Score & 0.721 & 0.794 \\
Fast-DetectGPT & 0.627 & 0.433 \\
Binoculars & 0.579 & 0.199 \\
\bottomrule
\end{tabular}
\caption{LOGO macro F1 for ordinal depth classification. 7-feature ensemble significantly outperforms all baselines (bootstrap $p < 0.01$). CI for avg-logprob is wide: [0.616, 0.820], reflecting high cross-generator variance.}
\label{tab:baseline}
\end{table}

The 7-feature ensemble achieves 0.798, significantly outperforming the best single detector (avg-logprob 0.765; bootstrap $\Delta = 3.3$ pp, $p < 0.01$).
Avg-logprob is strongest among single detectors because it directly measures NLL mean --- essentially an unnormalized entropy mean.
However, its narrow 95\% CI [0.616, 0.820] reveals high per-generator variance, whereas the 7-feature ensemble achieves a much tighter CI [0.786, 0.799].

Binoculars performs worst (0.579): its two-model probability ratio partially cancels depth signal, as both models' NLL decreases with depth.
Fast-DetectGPT (0.627) was designed for binary detection; its perturbation-based curvature metric is insensitive to ordinal depth variation.

\subsection{Cross-Model Chains}
\label{sec:crossmodel}

Real-world contamination involves diverse generators. We evaluate on cross-model chains where each depth step uses a different generator (Table~\ref{tab:crossmodel}).

\begin{table}[t]
\centering
\small
\begin{tabular}{lcc}
\toprule
\textbf{Setting} & \textbf{Binary F1} & \textbf{3-class F1} \\
\midrule
Single-model (LOGO) & 0.694 & --- \\
Cross-model (LOGO) & \textbf{0.830} & 0.789 \\
\bottomrule
\end{tabular}
\caption{Cross-model chains (5-gen, generator-level LOGO) vs.\ single-model LOGO. Generator heterogeneity \emph{amplifies} depth signal ($+$13.6 pp binary F1).}
\label{tab:crossmodel}
\end{table}

Cross-model binary F1 = 0.830 (95\% CI [0.812, 0.848]) exceeds single-model LOGO by 13.6 pp.
Per-generator variance drops to 4.0 pp (vs.\ 5.4 pp for single-model), indicating more stable signal under heterogeneity.
DDR drifting-feature correlations strengthen in cross-model chains (hapax $\rho = -0.72$ vs.\ $-0.647$ single-model): each generator change injects fresh distributional perturbation, preventing convergence to a model-specific attractor state.
This is a positive ecological validity result: realistic multi-generator contamination is \emph{easier}, not harder, to detect.

\subsection{Downstream Utility (Negative Finding)}
\label{sec:downstream}

We test whether depth-aware data curation improves downstream language model quality by fine-tuning OLMo-2-1B (1.485B parameters) on depth-stratified subsets (Table~\ref{tab:downstream}).

\begin{table}[t]
\centering
\small
\begin{tabular}{lrcc}
\toprule
\textbf{Variant} & $n$ & \textbf{PPL} & \textbf{$\pm$ std} \\
\midrule
All data & 12,299 & 2.335 & 0.0003 \\
Size-matched (random) & 3,951 & 2.367 & 0.0002 \\
Depth 0--1 only & 3,951 & 2.387 & 0.001 \\
Binary-filtered & 3,951 & 2.389 & 0.001 \\
Depth-0 only & 797 & 2.681 & 0.001 \\
\bottomrule
\end{tabular}
\caption{Downstream CLM perplexity (3 seeds). Depth-aware curation (depth 0--1, binary-filtered) does not outperform random subsampling at equal sample count.}
\label{tab:downstream}
\end{table}

The ranking is: all data (2.335) $<$ size-matched random (2.367) $<$ depth 0--1 (2.387) $\approx$ binary-filtered (2.389) $\ll$ depth-0 only (2.681).
TOST equivalence tests at $\pm 2\%$ margin confirm that depth-aware curation and random subsampling produce equivalent downstream quality ($p = 1.1 \times 10^{-4}$ for size-matched vs.\ depth 0--1; $p = 9 \times 10^{-5}$ for size-matched vs.\ binary-filtered).
At the stricter $\pm 1\%$ margin, equivalence is not confirmed ($p = 0.08$).

Depth filtering removes deeper text but also removes diverse training signal; random subsampling preserves more complete distributional coverage.
Seed variance is minimal ($< 0.002$ across seeds), confirming the stability of this negative finding.
This establishes depth fingerprinting as a \textbf{diagnostic} tool --- useful for characterizing data provenance --- rather than a curative tool for improving training data quality.

%=============================================================================
% 5. DISCUSSION
%=============================================================================
\section{Discussion}

\paragraph{DDR as location shift.}
The drifting/stationary split is consistent with recursive generation inducing a location shift in the NLL distribution: central-tendency features (mean, diversity) drift while shape features (tail decay, entropy max) are preserved.
This is not a ``discovery'' in the strong sense --- it follows from properties of the NLL distribution under recursive compression --- but its empirical robustness across 8 generators, 5 families, and 3 reference models is the contribution.

\paragraph{Entropy std as intermediate case.}
Entropy std shows weak drift ($|\rho| \approx 0.20$, $R^2 = 0.41$), with magnitude $\sim$1/3 of the drifting group.
We classify it as ``weakly drifting'' without disrupting the core binary distinction.
Its feature importance is low (rank 5/7 in the 5-gen evaluation).

\paragraph{Limitations.}
\begin{enumerate}[nosep,leftmargin=*]
\item \textbf{Base models only.} Instruction tuning compresses the distributional variation DDR exploits (LOGO F1 = 0.338 for OLMo-Instruct vs.\ 0.763 base). This is the primary scope limitation.
\item \textbf{Temperature sensitivity.} Reliable within $t \in [0.5, 0.9]$. At $t = 1.0$, entropy mean $\rho$ drops to $-0.002$.
\item \textbf{Scale effect.} Larger generators show moderate degradation (Qwen3-8B binary F1 = 0.683; $\sim$85\% of drifting $|\rho|$ vs.\ $<$3B average). DDR pattern is qualitatively preserved (\S\ref{sec:scale}).
\item \textbf{3-class resolution ceiling.} 5-class F1 drops to $\sim$0.50; ordinal methods offer no improvement (\S\ref{sec:ordinal}).
\item \textbf{Prompt-type dependence.} Continuation chains (F1 = 0.781) outperform elaboration (0.730) by 5.1 pp (\S\ref{sec:prompttype}).
\item \textbf{Downstream evidence.} Limited to single-task (CLM) and single-metric (perplexity) evaluation with 3 seeds.
\end{enumerate}

%=============================================================================
% 6. CONCLUSION
%=============================================================================
\section{Conclusion}

We formulate generation depth classification and characterize differential depth response (DDR) across 8 base-model generators from 5 families.
A 7-feature distributional ensemble achieves leak-free binary F1 = 0.694 (per-generator range: 0.633--0.724) under LOGO evaluation, with cross-model chains reaching 0.830.
DDR is robust across reference models and generator families but degrades under instruction tuning, defining a clear scope boundary.
Depth-aware curation is statistically equivalent to random subsampling for downstream quality, positioning depth fingerprinting as a diagnostic rather than curative tool.
Future work should address instruction-tuned models, larger scales ($>$14B), and downstream tasks beyond causal language modeling.

%=============================================================================
% REFERENCES
%=============================================================================
{\small
\bibliographystyle{plainnat}
\begin{thebibliography}{20}

\bibitem[Bao et~al.(2024)]{bao2024fast}
Bao, G., Zhao, Y., Teng, Z., Yang, L., and Zhang, Y.
\newblock Fast-DetectGPT: Efficient zero-shot detection of machine-generated text via conditional probability curvature.
\newblock In \emph{Proc.\ ICLR}, 2024.

\bibitem[Briakou et~al.(2025)]{briakou2025llm}
Briakou, E., et~al.
\newblock LLM as a broken telephone: Iterative paraphrasing converges to attractor cycles.
\newblock In \emph{Proc.\ ACL}, 2025.

\bibitem[Dohmatob et~al.(2024)]{dohmatob2024tale}
Dohmatob, E., Feng, Y., and Gu, Q.
\newblock A tale of tails: Model collapse as a change of scaling laws.
\newblock In \emph{Proc.\ ICML}, 2024.

\bibitem[Gehrmann et~al.(2019)]{gehrmann2019gltr}
Gehrmann, S., Strobelt, H., and Rush, A.~M.
\newblock GLTR: Statistical detection and visualization of generated text.
\newblock In \emph{Proc.\ ACL Demo}, 2019.

\bibitem[Hans et~al.(2024)]{hans2024spotting}
Hans, A., Schwarzschild, A., Cheber, V., Alon, U., Lyu, Z., and Goldstein, T.
\newblock Spotting LLMs with Binoculars: Zero-shot detection of machine-generated text.
\newblock In \emph{Proc.\ ICML}, 2024.

\bibitem[Liang et~al.(2024)]{liang2024monitoring}
Liang, W., et~al.
\newblock Monitoring AI-modified content at scale: A case study on the impact of ChatGPT on AI conference peer reviews.
\newblock \emph{arXiv:2403.07183}, 2024.

\bibitem[McCarthy and Jarvis(2010)]{mccarthy2010mtld}
McCarthy, P.~M. and Jarvis, S.
\newblock MTLD, vocd-D, and HD-D: A validation study of sophisticated approaches to lexical diversity assessment.
\newblock \emph{Behavior Research Methods}, 42(2):381--392, 2010.

\bibitem[Mitchell et~al.(2023)]{mitchell2023detectgpt}
Mitchell, E., Lee, Y., Khazatsky, A., Manning, C.~D., and Finn, C.
\newblock DetectGPT: Zero-shot machine-generated text detection using probability curvature.
\newblock In \emph{Proc.\ ICML}, 2023.

\bibitem[Neal and Woodard(2017)]{neal2017surveying}
Neal, T. and Woodard, D.
\newblock Surveying stylometry techniques and applications.
\newblock \emph{ACM Computing Surveys}, 50(6):1--36, 2017.

\bibitem[Shumailov et~al.(2024)]{shumailov2024model}
Shumailov, I., Shumaylov, Z., Zhao, Y., Gal, Y., Papernot, N., and Anderson, R.
\newblock AI models collapse when trained on recursively generated data.
\newblock \emph{Nature}, 631:755--759, 2024.

\bibitem[Tulchinskii et~al.(2024)]{tulchinskii2024intrinsic}
Tulchinskii, E., et~al.
\newblock Intrinsic dimension estimation for robust detection of AI-generated texts.
\newblock In \emph{Proc.\ NeurIPS}, 2024.

\bibitem[Wenzek et~al.(2020)]{wenzek2020ccnet}
Wenzek, G., Lachaux, M.-A., Conneau, A., Chaudhary, V., Guzm{\'a}n, F., Joulin, A., and Grave, E.
\newblock CCNet: Extracting high quality monolingual datasets from web crawl data.
\newblock In \emph{Proc.\ LREC}, 2020.

\bibitem[Yang et~al.(2024)]{yang2024dna}
Yang, X., Cheng, W., Petzold, L., Wang, W.~Y., and Chen, H.
\newblock DNA-GPT: Divergent n-gram analysis for training-free detection of GPT-generated text.
\newblock In \emph{Proc.\ ICLR}, 2024.

\end{thebibliography}
}

%=============================================================================
% APPENDIX
%=============================================================================
\appendix
\onecolumn

\section{Prompt-Type Comparison}
\label{sec:prompttype}

We evaluate three prompt types --- continuation, elaboration, and inspired rewriting --- on Pythia-2.8B with 1,000 seeds (Table~\ref{tab:prompttype}).

\begin{table}[h]
\centering
\small
\begin{tabular}{lccc}
\toprule
\textbf{Prompt Type} & \textbf{CV F1} & \textbf{LOPTO F1} & \textbf{Retention (d4/d0)} \\
\midrule
Continuation & 0.781 & 0.762 & 0.831 \\
Inspired rewriting & 0.737 & 0.733 & 0.939 \\
Elaboration & 0.730 & 0.762 & 0.986 \\
\bottomrule
\end{tabular}
\caption{Per-prompt-type classification (CV = within-type, LOPTO = leave-one-prompt-type-out). Continuation is strongest (5.1 pp gap) because unconstrained generation maximizes distributional drift. LOPTO average 0.753 confirms cross-prompt-type deployment is viable.}
\label{tab:prompttype}
\end{table}

DDR is qualitatively consistent across all prompt types: all drifting features show $|\rho| > 0.55$.
Continuation produces the strongest signal because it maximizes recursive compression without semantic constraints.
Elaboration is weakest because semantic-preservation constraints limit distributional drift amplitude.
The depth retention difference (elaboration 98.6\% vs.\ continuation 83.1\%) reflects that continuation's deeper outputs are more likely to be truncated by minimum-length filters.

Feature drift magnitudes by prompt type:

\begin{table}[h]
\centering
\small
\begin{tabular}{lccc}
\toprule
\textbf{Feature} & \textbf{Cont.\ $\rho$} & \textbf{Elab.\ $\rho$} & \textbf{Rewrite $\rho$} \\
\midrule
Hapax ratio & $-0.651$ & $-0.674$ & $-0.670$ \\
TTR & $-0.629$ & $-0.666$ & $-0.659$ \\
Entropy mean & $-0.665$ & $-0.583$ & $-0.564$ \\
\bottomrule
\end{tabular}
\caption{Spearman $\rho$ by feature and prompt type. DDR pattern is preserved across all types.}
\end{table}


\section{Paraphrastic vs.\ Non-Paraphrastic Chains}
\label{sec:paraphrastic}

Paraphrastic chains (each step paraphrases the previous output) exhibit substantially weaker depth signal than continuation chains (Table~\ref{tab:para}).

\begin{table}[h]
\centering
\small
\begin{tabular}{lcc}
\toprule
\textbf{Chain Type} & \textbf{Binary F1} & \textbf{3-class F1} \\
\midrule
Continuation & 0.781 & 0.834 \\
Paraphrastic & 0.618 & 0.680 \\
\midrule
$\Delta$ & $-16.3$ pp & $-15.4$ pp \\
\bottomrule
\end{tabular}
\caption{Paraphrastic vs.\ continuation chains. The 16.3 pp binary gap reflects attractor convergence in paraphrastic chains.}
\label{tab:para}
\end{table}

Paraphrastic chains converge by depth $\sim$3 (depth-3 vs.\ depth-2 binary F1 = 0.525, near chance), consistent with attractor cycle theory \citep{briakou2025llm}.
Continuation chains retain discriminability at deeper depths (depth-3 vs.\ depth-2 F1 = 0.640).
Semantic similarity from depth 0 to depth 4 is 0.88 for paraphrastic chains (meaning preserved) vs.\ 0.52 for continuation (significant topic drift).
DDR correlations in paraphrastic chains are uniformly weak ($|\rho| < 0.33$), confirming that DDR is specific to non-paraphrastic generation.


\section{Multi-Domain Evaluation}
\label{sec:multidomain}

We evaluate on four seed domains to test ecological validity (Table~\ref{tab:domain}).

\begin{table}[h]
\centering
\small
\begin{tabular}{lcccc}
\toprule
\textbf{Domain} & \textbf{Binary F1} & \textbf{3-class F1} & \textbf{Hapax $\rho$} & \textbf{Ent.\ mean $\rho$} \\
\midrule
Wikipedia & 0.694 & 0.770 & $-0.647$ & $-0.610$ \\
News & 0.712 & 0.715 & $-0.621$ & $-0.595$ \\
Code & 0.687 & 0.673 & $-0.630$ & $-0.581$ \\
Academic & 0.628 & 0.656 & $-0.543$ & $-0.520$ \\
\midrule
Cross-domain & 0.676 & --- & --- & --- \\
\bottomrule
\end{tabular}
\caption{Per-domain evaluation. DDR pattern is qualitatively consistent (all drifting $|\rho| > 0.52$). Academic domain is weakest due to formulaic genre conventions limiting distributional drift.}
\label{tab:domain}
\end{table}

Maximum domain gap is 8.4 pp (news 0.712 vs.\ academic 0.628).
Academic text is weakest because its highly formulaic genre conventions (strict sentence structures, specialized terminology) limit the distributional drift space, and its low baseline TTR creates a ceiling effect.
News text is strongest due to high vocabulary diversity providing larger feature margins.
All domains show drifting $|\rho| > 0.52$, confirming DDR pattern consistency.


\section{Finer-Grained Ordinal Classification}
\label{sec:ordinal}

Five-class ordinal classification (gen-0 through gen-4) reveals a practical 3-class resolution ceiling (Table~\ref{tab:ordinal}).

\begin{table}[h]
\centering
\small
\begin{tabular}{lcc}
\toprule
\textbf{Method} & \textbf{5-class F1} & \textbf{3-class F1} \\
\midrule
XGBoost & 0.503 & 0.773 \\
CORAL & 0.497 & --- \\
Frank-Hall & 0.501 & --- \\
\bottomrule
\end{tabular}
\caption{5-class ordinal classification. All methods perform comparably ($<$1 pp spread). 3-class reference included for comparison.}
\label{tab:ordinal}
\end{table}

Per-class 5-class F1: gen-0 = 0.952, gen-1 = 0.518, gen-2 = 0.348, gen-3 = 0.325, gen-4 = 0.372.
Gen-0 remains highly discriminable (human text); gen-2/3/4 are severely confused (mutual misclassification rate 42\%).
The bottleneck is the gen-1/gen-2 boundary, which concentrates 86\% of binary errors as false negatives.
Ordinal methods (CORAL, Frank-Hall) do not improve over standard XGBoost because the ordinal margin between adjacent deeper classes is not uniform --- the gen-0$\to$gen-1 margin far exceeds gen-2$\to$gen-3, violating the equal-spacing assumption these methods require.
The MAE of 0.72 depth levels indicates predictions are off by less than one class on average.


\section{Human Edit Robustness}
\label{sec:edit}

We test robustness to simulated human editing via random word replacement and LLM-based rephrasing (Table~\ref{tab:edit}).

\begin{table}[h]
\centering
\small
\begin{tabular}{lcc}
\toprule
\textbf{Condition} & \textbf{Binary F1} & \textbf{$\Delta$ (pp)} \\
\midrule
Baseline (no edit) & 0.695 & --- \\
5\% random replacement & 0.650 & $-4.6$ \\
10\% random replacement & 0.614 & $-8.1$ \\
20\% random replacement & 0.591 & $-10.4$ \\
LLM rephrase & 0.647 & $-4.5$ \\
\bottomrule
\end{tabular}
\caption{Binary F1 under simulated editing. Degradation is graceful: LLM rephrase ($\sim$50\% edit distance, $\sim$1.7 sentences replaced) causes only 4.5 pp drop.}
\label{tab:edit}
\end{table}

Degradation is graceful rather than catastrophic.
At 5\% random replacement, distributional features (TTR, entropy, hapax) are barely affected because they are global statistics --- sparse point noise does not alter aggregate patterns.
LLM rephrasing ($\sim$50\% edit distance) causes a comparable 4.5 pp drop: rephrasing preserves coherence and topic but alters local word choice, partially disrupting surface features while reference features (NLL-based) remain more robust because rephrased text's NLL retains depth-dependent patterns.
Even under 20\% replacement, binary F1 (0.591) substantially exceeds chance (0.50).
Multiclass degradation is larger ($-7.1$ pp under LLM rephrase) because editing primarily blurs the gen-1/gen-2+ boundary.


\section{LD-Score as Additional Feature}
\label{sec:ldscore}

We evaluate LD-Score (character-level distribution signature) as a potential 8th feature (Table~\ref{tab:ldscore}).

\begin{table}[h]
\centering
\small
\begin{tabular}{lc}
\toprule
\textbf{Configuration} & \textbf{LOGO F1} \\
\midrule
7-feature ensemble & 0.798 \\
8-feature (+ LD-Score) & 0.801 \\
LD-Score standalone & 0.721 \\
\bottomrule
\end{tabular}
\caption{Adding LD-Score to the 7-feature ensemble yields only 0.3 pp improvement. LD-Score correlates heavily with surface features (hapax $r = 0.72$, TTR $r = 0.68$).}
\label{tab:ldscore}
\end{table}

LD-Score alone achieves moderate signal (F1 = 0.721, Spearman $\rho = 0.794$ with depth), but adding it to the ensemble improves F1 by only 0.3 pp.
Character-level vocabulary diversity is highly co-linear with word-level surface features (hapax $r = 0.72$, TTR $r = 0.68$), providing negligible complementary information.
We retain the 7-feature set as the recommended configuration for simplicity, interpretability, and efficiency.


\section{End-to-End Neural Classifier}
\label{sec:qlora}

We compare the 7-feature ensemble against QLoRA-finetuned end-to-end classifiers as an upper-bound reference (Table~\ref{tab:qlora}).

\begin{table}[h]
\centering
\small
\begin{tabular}{lccc}
\toprule
\textbf{Method} & \textbf{3-class F1} & \textbf{Binary F1} & \textbf{GPU (GB)} \\
\midrule
QLoRA (instruct) & 0.897 & 0.852 & 8.2 \\
QLoRA (base) & 0.881 & 0.826 & 8.2 \\
7-feature ensemble & 0.773 & 0.694 & 2.4 \\
\bottomrule
\end{tabular}
\caption{QLoRA vs.\ 7-feature ensemble. QLoRA achieves substantially higher within-distribution F1 but at 3.4$\times$ GPU memory and 6.5$\times$ inference time.}
\label{tab:qlora}
\end{table}

QLoRA achieves 12--16 pp higher within-distribution F1 by learning contextual, positional, and discourse-level patterns beyond our handcrafted features.
However, QLoRA's cross-generator generalization (LOGO binary 0.720) is substantially below its within-distribution performance (0.826--0.852), indicating partial overfitting to generator-specific token patterns.
The 7-feature ensemble's advantages are: (a) 3.4$\times$ less GPU memory, (b) 6.5$\times$ faster inference, (c) interpretable features with clear DDR semantics, (d) more stable cross-generator generalization.
Notably, QLoRA's 5-class F1 (0.551) only exceeds the handcrafted baseline (0.503) by $\sim$5 pp, confirming that finer classification is bottlenecked by task-inherent difficulty.


\section{Text Length Confound Control}
\label{sec:length}

We verify that depth signal is not an artifact of text length variation (Table~\ref{tab:length}).

\begin{table}[h]
\centering
\small
\begin{tabular}{lcc}
\toprule
\textbf{Condition} & \textbf{Binary F1} & \textbf{$\Delta$ (pp)} \\
\midrule
Unmatched & 0.770 & --- \\
Random subsample (equal $n$) & 0.741 & $-2.9$ \\
Length-matched & 0.724 & $-4.6$ \\
\bottomrule
\end{tabular}
\caption{Length confound analysis. Of the 4.6 pp drop from length matching, 2.9 pp is attributable to sample-size reduction and only 1.7 pp to deconfounding.}
\label{tab:length}
\end{table}

Length--depth Spearman $\rho$ is weak ($[-0.066, +0.134]$): text length does not systematically vary with depth.
Of the 4.6 pp drop from length matching, 2.9 pp is attributable to sample-size reduction (random subsample control) and only 1.7 pp to actual deconfounding.
Length-normalized surface features (computed over fixed-length windows) \emph{strengthen} the depth trend ($|\rho|$ increases), confirming that DDR arises from intrinsic vocabulary diversity decrease rather than length variation.


\section{Large-Scale Generator}
\label{sec:scale}

We add Qwen3-8B-Base to test scale effects (Table~\ref{tab:scale}).

\begin{table}[h]
\centering
\small
\begin{tabular}{lccc}
\toprule
\textbf{Generator} & \textbf{Binary F1} & \textbf{Hapax $\rho$} & \textbf{Tail decay $\rho$} \\
\midrule
$<$3B average & 0.700 & $-0.647$ & $-0.040$ \\
Qwen3-8B & 0.683 & $-0.558$ & $-0.089$ \\
\bottomrule
\end{tabular}
\caption{Scale effect: 8B generator shows moderate degradation ($-1.7$ pp binary F1) with attenuated drifting $|\rho|$ ($\sim$85\% of $<$3B average). DDR pattern (drifting vs.\ stationary split) is qualitatively preserved.}
\label{tab:scale}
\end{table}

Qwen3-8B binary F1 = 0.683 (95\% CI [0.660, 0.706]), a moderate 1.7 pp drop from the $<$3B average.
Drifting features retain $\sim$85\% of their correlation magnitude (hapax $\rho = -0.558$ vs.\ $-0.647$).
The DDR pattern is qualitatively preserved: the three drifting features remain highest-$|\rho|$ (hapax $-0.558$, TTR $-0.542$, entropy mean $-0.475$), while stationary features remain low ($|\rho| < 0.10$).
Higher-capacity models produce more concentrated generation distributions with lower baseline vocabulary diversity, compressing depth-wise feature margins.
$N = 4{,}497$ total records with depth-dependent dropout (d4: 755/995 seeds) more pronounced than for smaller models.

\end{document}
